%% !TEX root = manuscript.tex

\chapter{Reporting et missions SI transverses}
\minitoc
\newpage

Pendant que les deux reprises du chapitre précédent se déroulaient, j'ai maintenu mes
activités de gestionnaire d'information : les demandes de reporting des services et une
contribution à l'infrastructure système. Conduites sur les deux semestres, elles sont exposées
en détail dans mon mémoire de mi-parcours~; je n'en retiens ici que ce qui éclaire la suite,
car c'est en les menant que j'ai rencontré, sous une forme miniature, le problème qui commande
le projet de Data Science du chapitre~3 : lorsqu'une règle de gestion n'a nulle part où être
écrite, elle finit par être recalculée partout, différemment.

\hypertarget{le-maintien-opuxe9rationnel-du-reporting}{%
\section{Le maintien opérationnel du reporting}\label{le-maintien-opuxe9rationnel-du-reporting}}

Ces demandes consistent à traduire des besoins métier parfois abstraits en indicateurs fiables.
Le travail commence rarement par une spécification : il commence par une gêne exprimée en
réunion, ou par un tableau que quelqu'un refait à la main tous les mois. Reformuler la demande
jusqu'à ce qu'elle soit calculable --- de quelles lignes parle-t-on, sur quelle période, que
fait-on des cas limites~? --- est la partie que je sous-estimais en arrivant, et c'est celle qui
décide de l'utilité du résultat.

La demande du service juridique en donne l'exemple le plus significatif. Le besoin portait sur
le suivi des commandes clients, avec une règle de gestion en alternance stricte : la première
commande est payante, la deuxième gratuite, la troisième payante, et ainsi de suite. La
difficulté tient à ce que cette règle ne porte pas sur une ligne mais sur la position d'une
ligne dans une séquence~: aucune colonne de la base ne contient l'information «~c'est la
deuxième commande de ce client sur ce lot~», qui n'existe que si on la calcule et change dès
qu'une commande est ajoutée ou annulée.

Deux contraintes interdisaient de stocker le résultat : je ne possédais pas les droits
d'écriture permettant de créer un champ de statut dans la base SQL, et les fichiers Excel
sources étaient écrasés à chaque mise à jour. Il fallait donc un calcul dynamique, exécuté à
la génération du tableau : identifier les opérations et les lots, classer les commandes par
ordre chronologique, attribuer un rang, puis appliquer la règle --- rang impair, statut
«~Payant~»~; rang pair, «~Gratuit~». La Figure 2.1 représente cette chaîne de calcul.

\begin{figure}[!htbp]
\centering
\includegraphics[width=0.65\textwidth]{fig_2_1_regle_facturation.png}
\caption{La règle de facturation juridique traduite en calcul de rang, exécuté à chaque génération du tableau}
\end{figure}

Cette approche répond au besoin sans altérer les données sources, mais elle a un coût : la
règle est recalculée à chaque affichage et n'existe nulle part ailleurs que dans la définition
du tableau. Si un autre service a besoin de la même notion, rien ne garantit qu'elle sera
implantée à l'identique. C'est exactement le problème que le projet du chapitre~3 résoudra, en
écrivant les définitions partagées une seule fois dans des vues de base de données. Le régime
de preuve retenu ici est celui de la \textbf{recette par l'usage}~: seul le service demandeur
pouvait juger de la conformité de la règle appliquée, et cette validation ne dit rien, à court
terme, de l'utilité de l'outil dans la durée.

\hypertarget{le-duxe9ploiement-du-single-sign-on}{%
\section{Le déploiement du Single Sign-On}\label{le-duxe9ploiement-du-single-sign-on}}

J'ai par ailleurs collaboré avec l'administratrice des systèmes et outils SI au déploiement du
SSO (\emph{Single Sign-On}) sur les postes de deux périmètres, l'agence de Perpignan et le service
financier du siège : installation du connecteur, vérification de la synchronisation avec
l'annuaire d'entreprise, accompagnement des utilisateurs à leur première connexion. Nous avons
volontairement commencé par ces deux périmètres avant d'envisager une généralisation.

Cette mission m'a appris quelque chose que je n'attendais pas d'un déploiement. L'annuaire
d'entreprise sur lequel repose le SSO est un référentiel d'identités : il répond, pour les
collaborateurs, à la question que je me posais pour les clients au chapitre~1, à savoir
comment garantir qu'une personne n'existe qu'une fois et qu'on sait la reconnaître. Voir ce
problème résolu du côté de l'infrastructure éclaire par contraste la situation côté données de
gestion, où aucun référentiel de ce type n'existait avant la reprise des tiers.
L'accompagnement des utilisateurs a été instructif à sa manière : une solution techniquement
correcte qui déroute l'utilisateur au moment de s'en servir n'est pas une solution livrée.

\hypertarget{ce-que-les-missions-filuxe9es-couxfbtent-et-ce-quelles-rapportent}{%
\section{Ce que les missions filées coûtent et ce qu'elles rapportent}\label{ce-que-les-missions-filuxe9es-couxfbtent-et-ce-quelles-rapportent}}

Ces missions transverses fragmentent le temps, et un chantier de reprise interrompu trois fois
dans la journée pour une demande de tableau de bord se reprend mal, puisqu'il faut à chaque
fois reconstituer où l'on en était dans une logique de formules imbriquées. Cette
fragmentation a été la difficulté la plus constante de l'année, et ce que j'ai fini par mettre
en place tient en peu de choses : réserver les plages longues aux travaux qui ne se découpent
pas, regrouper les demandes courtes, et surtout écrire où j'en étais avant d'interrompre une
tâche.

Elles rapportent cependant trois choses que je n'aurais pas obtenues autrement : une
\textbf{connaissance du métier} que la fréquentation des bases ne donne pas, car on ne comprend pas
ce qu'est une commande gratuite en lisant une table~; une \textbf{crédibilité} auprès des services,
qui a rendu possible au second semestre un échange direct avec la direction financière~; et
une \textbf{vision de l'architecture} du système d'information, depuis l'authentification des postes
jusqu'aux tableaux de bord, qui m'a permis de concevoir le projet du second semestre comme une
chaîne complète plutôt que comme une suite de modèles.

Ces solutions restent des contournements : une règle recalculée à chaque affichage tient tant
que personne d'autre n'en a besoin. Elles auront néanmoins maintenu le système décisionnel en
état de marche pendant les douze mois qu'a duré la reconstruction du socle. Ce socle étant
désormais fiable, le chapitre suivant expose ce qu'il devient possible d'en faire.
